\subsection{The Six FEPSDE Dimensions}
\label{subsec:fepsde-dimensions}

%%%%%%%%%%%%%%%%%%%%%%%%%%%%%%%%%%%%%%%%%%%%%%%%%%%%%%%%%%%%%%%%%%%%%%%%%%%%%%%
% 10.2: The six welfare dimensions that apply to ALL utility categories
%%%%%%%%%%%%%%%%%%%%%%%%%%%%%%%%%%%%%%%%%%%%%%%%%%%%%%%%%%%%%%%%%%%%%%%%%%%%%%%

Section~\ref{subsec:three-categories} established \emph{for whom} utility matters 
(INU, KNU, IDN). This section establishes \emph{in what respects} utility matters:
the six FEPSDE dimensions.

\subsubsection{Why Multiple Dimensions?}

\paragraph{The Inadequacy of Single-Dimensional Welfare.}
Traditional welfare economics uses a single metric: income, consumption, or GDP.
This is convenient but inadequate:

\begin{itemize}
\item \textbf{The Easterlin Paradox:} Above ~\$20,000 GDP per capita, 
more income does not increase reported happiness \citep{easterlin1974does}

\item \textbf{The American Paradox:} The United States has high GDP but 
lower life expectancy, higher inequality, and lower life satisfaction 
than peer countries

\item \textbf{The Mismeasurement Problem:} GDP counts pollution cleanup 
as positive, ignores household production, and treats depletion of 
natural capital as income
\end{itemize}

\paragraph{The FEPSDE Solution.}
We propose six dimensions that together capture the full spectrum of 
human welfare. Each dimension is:

\begin{enumerate}
\item \textbf{Empirically distinct:} Captures variance not captured by others
\item \textbf{Universally relevant:} Applies across cultures, contexts, time
\item \textbf{Measurable:} Can be operationalized for empirical work
\item \textbf{Policy-relevant:} Corresponds to distinct policy domains
\item \textbf{Context-sensitive:} Responds to $\Psi$-dimensions via $\Gamma$-matrix
\end{enumerate}

%%%%%%%%%%%%%%%%%%%%%%%%%%%%%%%%%%%%%%%%%%%%%%%%%%%%%%%%%%%%%%%%%%%%%%%%%%%%%%%
\subsubsection{The Six Dimensions Defined}
%%%%%%%%%%%%%%%%%%%%%%%%%%%%%%%%%%%%%%%%%%%%%%%%%%%%%%%%%%%%%%%%%%%%%%%%%%%%%%%

\begin{definition}[FEPSDE Welfare Dimensions]
\label{def:fepsde-dimensions}
\begin{align*}
U^F &: \text{\textbf{Financial} --- material wealth, income, economic security} \\[0.3em]
U^E &: \text{\textbf{Emotional} --- psychological wellbeing, satisfaction, meaning} \\[0.3em]
U^P &: \text{\textbf{Physical} --- health, energy, bodily integrity, longevity} \\[0.3em]
U^S &: \text{\textbf{Social} --- relationships, belonging, trust, status} \\[0.3em]
U^D &: \text{\textbf{Digital} --- connectivity, access, digital participation} \\[0.3em]
U^{Eco} &: \text{\textbf{Ecological} --- environmental quality, sustainability}
\end{align*}
\end{definition}

%%%%%%%%%%%%%%%%%%%%%%%%%%%%%%%%%%%%%%%%%%%%%%%%%%%%%%%%%%%%%%%%%%%%%%%%%%%%%%%
\subsubsection{Financial Dimension ($U^F$)}
%%%%%%%%%%%%%%%%%%%%%%%%%%%%%%%%%%%%%%%%%%%%%%%%%%%%%%%%%%%%%%%%%%%%%%%%%%%%%%%

\paragraph{What It Captures.}
Material resources and economic security:
\begin{itemize}
\item Income (wages, salaries, transfers)
\item Wealth (assets minus liabilities)
\item Consumption possibilities
\item Economic security (protection against shocks)
\item Access to credit and financial services
\end{itemize}

\paragraph{Why It's Necessary But Not Sufficient.}
Financial resources are instrumentally valuable---they enable consumption 
of goods and services. But the relationship between money and welfare 
is complex:

\begin{itemize}
\item \textbf{Diminishing returns:} The marginal utility of income decreases
\item \textbf{Relative position:} Satisfaction depends on comparison to others
\item \textbf{Adaptation:} People adapt to income changes (hedonic treadmill)
\item \textbf{Opportunity cost:} Earning money may sacrifice other dimensions
\end{itemize}

\paragraph{Measurement.}
\begin{itemize}
\item Income (individual, household, equivalized)
\item Wealth (net worth)
\item Consumption expenditure
\item Economic security indices
\item Gini coefficient (distribution)
\item Financial stress indicators
\end{itemize}

\paragraph{Context Sensitivity.}
From the $\Gamma$-matrix (Chapter~10.7):
\begin{itemize}
\item $\gamma_{F,E} = 0.42$: Economic development strongly predicts financial welfare
\item $\gamma_{F,I} = 0.31$: Institutions enable financial prosperity
\item $\gamma_{F,F} = 0.33$: Flexibility supports financial outcomes
\end{itemize}

%%%%%%%%%%%%%%%%%%%%%%%%%%%%%%%%%%%%%%%%%%%%%%%%%%%%%%%%%%%%%%%%%%%%%%%%%%%%%%%
\subsubsection{Emotional Dimension ($U^E$)}
%%%%%%%%%%%%%%%%%%%%%%%%%%%%%%%%%%%%%%%%%%%%%%%%%%%%%%%%%%%%%%%%%%%%%%%%%%%%%%%

\paragraph{What It Captures.}
Psychological wellbeing and subjective experience:
\begin{itemize}
\item Life satisfaction (cognitive evaluation)
\item Positive affect (happiness, joy, contentment)
\item Negative affect (anxiety, depression, stress)
\item Meaning and purpose (eudaimonic wellbeing)
\item Autonomy and self-determination
\item Mental health
\end{itemize}

\paragraph{Theoretical Foundation.}
\begin{itemize}
\item \textbf{Subjective Wellbeing:} \citet{diener1999subjective} established 
life satisfaction and affect as measurable constructs
\item \textbf{Self-Determination Theory:} \citet{ryan2000self} identified 
autonomy, competence, and relatedness as psychological needs
\item \textbf{PERMA Model:} \citet{seligman2011flourish} proposed Positive emotion, 
Engagement, Relationships, Meaning, Achievement
\item \textbf{Eudaimonia:} Aristotelian tradition emphasizes flourishing 
over hedonic pleasure
\end{itemize}

\paragraph{Why Emotional Welfare Is Not Reducible to Financial.}
\begin{itemize}
\item Nordic countries: High $U^E$ despite moderate post-tax income
\item Lottery winners: Temporary $U^E$ boost, then adaptation
\item Unemployment: $U^E$ drops far more than income loss would predict
\item Meaningful work: Often chosen over higher-paying alternatives
\end{itemize}

\paragraph{Measurement.}
\begin{itemize}
\item Life satisfaction scales (Cantril ladder, SWLS)
\item Positive/negative affect schedules (PANAS)
\item Depression and anxiety inventories (PHQ-9, GAD-7)
\item Meaning in life questionnaires (MLQ)
\item Experience sampling (momentary wellbeing)
\end{itemize}

\paragraph{Context Sensitivity.}
\begin{itemize}
\item $\gamma_{E,S} = 0.47$: Social trust is strongest predictor of emotional wellbeing
\item $\gamma_{E,T} = 0.29$: Temporal stability supports emotional welfare
\item $\gamma_{E,I} = 0.15$: Institutions matter but less than for financial
\end{itemize}

%%%%%%%%%%%%%%%%%%%%%%%%%%%%%%%%%%%%%%%%%%%%%%%%%%%%%%%%%%%%%%%%%%%%%%%%%%%%%%%
\subsubsection{Physical Dimension ($U^P$)}
%%%%%%%%%%%%%%%%%%%%%%%%%%%%%%%%%%%%%%%%%%%%%%%%%%%%%%%%%%%%%%%%%%%%%%%%%%%%%%%

\paragraph{What It Captures.}
Bodily health and functioning:
\begin{itemize}
\item Health status (absence of disease)
\item Physical functioning and mobility
\item Energy and vitality
\item Longevity and life expectancy
\item Freedom from pain
\item Sleep quality
\item Bodily integrity and safety
\end{itemize}

\paragraph{Why Health Is Fundamental.}
Health is both intrinsically valuable and instrumentally necessary:
\begin{itemize}
\item Severe illness reduces capacity to work ($U^F$)
\item Pain impairs emotional wellbeing ($U^E$)
\item Disability limits social participation ($U^S$)
\item Health problems dominate attention, crowding out other goals
\end{itemize}

As Schopenhauer observed: ``Health is not everything, but without health, 
everything is nothing.''

\paragraph{The Health-Wealth Paradox.}
\begin{itemize}
\item USA: Highest healthcare spending, lower life expectancy than peers
\item Costa Rica: Fraction of US income, similar life expectancy
\item Implication: $U^P$ is not simply purchased with $U^F$
\end{itemize}

\paragraph{Measurement.}
\begin{itemize}
\item Life expectancy and healthy life years (HALE)
\item Self-reported health status
\item Disability-adjusted life years (DALYs)
\item Quality-adjusted life years (QALYs)
\item Biomarkers (blood pressure, BMI, etc.)
\item Functional capacity assessments
\end{itemize}

\paragraph{Context Sensitivity.}
\begin{itemize}
\item $\gamma_{P,C} = 0.31$: Cognitive capacity (education) strongly predicts health
\item $\gamma_{P,I} = 0.28$: Institutions (healthcare systems) matter
\item $\gamma_{P,T} = 0.26$: Stability enables health investment
\end{itemize}

%%%%%%%%%%%%%%%%%%%%%%%%%%%%%%%%%%%%%%%%%%%%%%%%%%%%%%%%%%%%%%%%%%%%%%%%%%%%%%%
\subsubsection{Social Dimension ($U^S$)}
%%%%%%%%%%%%%%%%%%%%%%%%%%%%%%%%%%%%%%%%%%%%%%%%%%%%%%%%%%%%%%%%%%%%%%%%%%%%%%%

\paragraph{What It Captures.}
Relationships and social integration:
\begin{itemize}
\item Close relationships (family, friends, partners)
\item Social support networks
\item Community belonging
\item Social status and recognition
\item Trust (interpersonal and institutional)
\item Civic participation
\end{itemize}

\paragraph{Why Social Relationships Are Special.}
The empirical evidence is overwhelming:

\begin{itemize}
\item \textbf{Longevity:} Social relationships are the strongest predictor 
of longevity---stronger than smoking, obesity, or exercise \citep{holt2010social}

\item \textbf{Health:} Loneliness increases mortality risk by 26\%, 
comparable to smoking 15 cigarettes per day \citep{holt2015loneliness}

\item \textbf{Happiness:} Social relationships are the strongest predictor 
of life satisfaction across cultures \citep{helliwell2004social}

\item \textbf{Economics:} Social trust underpins economic transactions, 
reducing transaction costs \citep{knack1997does}
\end{itemize}

\paragraph{The Modern Social Crisis.}
\begin{itemize}
\item \textbf{Japan:} ``Hikikomori'' (social withdrawal), declining marriage, 
``kodokushi'' (lonely deaths)
\item \textbf{USA:} Decline in social capital since 1960s \citep{putnam2000bowling}
\item \textbf{Global:} Rising loneliness despite (or because of?) digital connectivity
\end{itemize}

\paragraph{Measurement.}
\begin{itemize}
\item Social network size and quality
\item Frequency of social contact
\item Loneliness scales (UCLA Loneliness Scale)
\item Trust surveys (World Values Survey)
\item Civic participation rates
\item Social support indices
\end{itemize}

\paragraph{Context Sensitivity.}
\begin{itemize}
\item $\gamma_{S,S} = 0.51$: Social trust context is strongest predictor of social welfare
\item $\gamma_{S,M} = -0.18$: Market integration can \emph{reduce} social welfare
\item $\gamma_{S,I} = 0.11$: Institutions have modest direct effect
\end{itemize}

The negative $\gamma_{S,M}$ is crucial: it captures the documented effect of 
trade shocks on community disruption \citep{autor2013china}.

%%%%%%%%%%%%%%%%%%%%%%%%%%%%%%%%%%%%%%%%%%%%%%%%%%%%%%%%%%%%%%%%%%%%%%%%%%%%%%%
\subsubsection{Digital Dimension ($U^D$)}
%%%%%%%%%%%%%%%%%%%%%%%%%%%%%%%%%%%%%%%%%%%%%%%%%%%%%%%%%%%%%%%%%%%%%%%%%%%%%%%

\paragraph{What It Captures.}
Participation in the digital world:
\begin{itemize}
\item Internet access and connectivity
\item Digital literacy and skills
\item Access to digital services (banking, government, healthcare)
\item Data rights and privacy
\item Protection from digital harms (cybercrime, misinformation)
\item Ability to disconnect (digital wellbeing)
\end{itemize}

\paragraph{Why Digital Is Now a Distinct Dimension.}
The COVID-19 pandemic revealed what was already becoming true:
digital access is now as fundamental as electricity or running water.

\begin{itemize}
\item \textbf{Economic participation:} Job applications, banking, commerce 
increasingly require digital access
\item \textbf{Social connection:} For many, primary mode of maintaining relationships
\item \textbf{Information access:} News, health information, education
\item \textbf{Government services:} Tax filing, benefit applications, identification
\end{itemize}

\paragraph{The Digital Divide.}
Those without digital access face compounding disadvantages:
\begin{itemize}
\item Cannot apply for many jobs ($U^F$ impact)
\item Cannot access telehealth ($U^P$ impact)
\item Cannot video-call family ($U^S$ impact)
\item Cannot participate in online discourse (civic participation)
\end{itemize}

\paragraph{The Digital Paradox.}
More is not always better:
\begin{itemize}
\item Heavy social media use correlates with lower $U^E$ \citep{twenge2018decreases}
\item ``Doomscrolling'' increases anxiety
\item Digital addiction is recognized clinical concern
\item Privacy erosion creates new vulnerabilities
\end{itemize}

\paragraph{Measurement.}
\begin{itemize}
\item Internet access (broadband, mobile)
\item Digital literacy assessments
\item E-government usage
\item Screen time and digital wellbeing indicators
\item Cybersecurity incident rates
\item Data privacy indices
\end{itemize}

\paragraph{Context Sensitivity.}
\begin{itemize}
\item $\gamma_{D,K} = 0.39$: Information access context predicts digital welfare
\item $\gamma_{D,C} = 0.27$: Cognitive capacity enables digital participation
\item $\gamma_{D,M} = 0.24$: Market scope expands digital opportunities
\end{itemize}

%%%%%%%%%%%%%%%%%%%%%%%%%%%%%%%%%%%%%%%%%%%%%%%%%%%%%%%%%%%%%%%%%%%%%%%%%%%%%%%
\subsubsection{Ecological Dimension ($U^{Eco}$)}
%%%%%%%%%%%%%%%%%%%%%%%%%%%%%%%%%%%%%%%%%%%%%%%%%%%%%%%%%%%%%%%%%%%%%%%%%%%%%%%

\paragraph{What It Captures.}
Environmental quality and sustainability:
\begin{itemize}
\item Air and water quality
\item Access to green spaces and nature
\item Climate stability
\item Biodiversity
\item Resource sustainability
\item Environmental justice (distribution of environmental harms)
\end{itemize}

\paragraph{Why Ecological Welfare Matters.}
\begin{itemize}
\item \textbf{Health impacts:} Air pollution causes 7 million premature deaths 
annually (WHO)
\item \textbf{Climate change:} Threatens all other dimensions through 
displacement, conflict, food insecurity
\item \textbf{Intergenerational equity:} Today's $U^F$ may come at cost of 
future generations' $U^{Eco}$
\item \textbf{Psychological benefits:} Nature exposure improves mental health ($U^E$)
\end{itemize}

\paragraph{The Growth-Environment Tension.}
The $\Gamma$-matrix reveals structural trade-offs:
\begin{itemize}
\item $\gamma_{Eco,E} = -0.19$: Economic development tends to reduce ecological welfare
\item $\gamma_{Eco,M} = -0.21$: Market integration tends to reduce ecological welfare
\end{itemize}

These negative coefficients are not estimation errors---they reflect 
the documented Environmental Kuznets Curve and globalization-environment tensions.

\paragraph{Measurement.}
\begin{itemize}
\item Air quality indices (PM2.5, ozone)
\item Water quality measures
\item Carbon footprint
\item Green space access
\item Ecological footprint
\item Environmental Performance Index
\end{itemize}

\paragraph{Context Sensitivity.}
\begin{itemize}
\item $\gamma_{Eco,T} = 0.23$: Stability enables long-term environmental investment
\item $\gamma_{Eco,I} = 0.21$: Institutions can enforce environmental protection
\item $\gamma_{Eco,E} = -0.19$: Economic development often harms environment
\item $\gamma_{Eco,M} = -0.21$: Globalization can export environmental damage
\end{itemize}

%%%%%%%%%%%%%%%%%%%%%%%%%%%%%%%%%%%%%%%%%%%%%%%%%%%%%%%%%%%%%%%%%%%%%%%%%%%%%%%
\subsubsection{Dimension Independence and Interaction}
%%%%%%%%%%%%%%%%%%%%%%%%%%%%%%%%%%%%%%%%%%%%%%%%%%%%%%%%%%%%%%%%%%%%%%%%%%%%%%%

\paragraph{Partial Independence.}
The six dimensions are partially independent---improvements in one 
do not automatically produce improvements in others:

\begin{table}[htbp]
\centering
\caption{Dimension Independence: Counterexamples}
\label{tab:dimension-independence}
\renewcommand{\arraystretch}{1.3}
\begin{tabular}{@{}lll@{}}
\toprule
\textbf{High in...} & \textbf{Can be low in...} & \textbf{Example} \\
\midrule
$U^F$ (Financial) & $U^E$ (Emotional) & Wealthy but depressed executive \\
$U^F$ (Financial) & $U^S$ (Social) & Rich but lonely \\
$U^P$ (Physical) & $U^E$ (Emotional) & Physically healthy, existential crisis \\
$U^D$ (Digital) & $U^S$ (Social) & Always online, no real friends \\
$U^F$ (Financial) & $U^{Eco}$ (Ecological) & Affluent in polluted city \\
\bottomrule
\end{tabular}
\end{table}

\paragraph{Systematic Interactions.}
Despite partial independence, dimensions do interact systematically:

\begin{itemize}
\item \textbf{F $\rightarrow$ P:} Income enables healthcare access
\item \textbf{S $\rightarrow$ E:} Social support buffers stress
\item \textbf{P $\rightarrow$ all:} Health is precondition for other pursuits
\item \textbf{F $\rightarrow$ Eco (negative):} Consumption drives environmental impact
\item \textbf{D $\rightarrow$ S (ambiguous):} Digital can enhance or substitute for social
\end{itemize}

These interactions are captured in the within-dimension correlations and 
in the $\Gamma$-matrix context sensitivities.

%%%%%%%%%%%%%%%%%%%%%%%%%%%%%%%%%%%%%%%%%%%%%%%%%%%%%%%%%%%%%%%%%%%%%%%%%%%%%%%
\subsubsection{FEPSDE Applied to Each Utility Category}
%%%%%%%%%%%%%%%%%%%%%%%%%%%%%%%%%%%%%%%%%%%%%%%%%%%%%%%%%%%%%%%%%%%%%%%%%%%%%%%

\paragraph{The Key Insight.}
Each FEPSDE dimension applies to \emph{each} utility category:

\begin{table}[htbp]
\centering
\caption{FEPSDE $\times$ Utility Categories: Examples}
\label{tab:fepsde-categories}
\small
\renewcommand{\arraystretch}{1.3}
\begin{tabular}{@{}llll@{}}
\toprule
\textbf{Dimension} & \textbf{INU} & \textbf{KNU} & \textbf{IDN} \\
\midrule
Financial & My income & Our family's wealth & ``I'm a provider'' \\
Emotional & My happiness & Our team morale & ``I'm a positive person'' \\
Physical & My health & Community health & ``I'm an athlete'' \\
Social & My friendships & Our cohesion & ``I'm a connector'' \\
Digital & My access & Our connectivity & ``I'm tech-savvy'' \\
Ecological & My environment & Our planet & ``I'm an environmentalist'' \\
\bottomrule
\end{tabular}
\end{table}

This creates the $3 \times 6 = 18$ base cells of the utility matrix,
which expand to 144 components when combined with time and valence
(Section~\ref{subsec:144-structure}).

%%%%%%%%%%%%%%%%%%%%%%%%%%%%%%%%%%%%%%%%%%%%%%%%%%%%%%%%%%%%%%%%%%%%%%%%%%%%%%%
\subsubsection{Summary: The Six FEPSDE Dimensions}
%%%%%%%%%%%%%%%%%%%%%%%%%%%%%%%%%%%%%%%%%%%%%%%%%%%%%%%%%%%%%%%%%%%%%%%%%%%%%%%

\begin{tcolorbox}[colback=gray!5!white, colframe=gray!75!black, title=Section 10.2 Summary]
\textbf{Six welfare dimensions:}
\begin{itemize}
\item \textbf{F (Financial):} Material resources, economic security
\item \textbf{E (Emotional):} Psychological wellbeing, meaning
\item \textbf{P (Physical):} Health, energy, longevity
\item \textbf{S (Social):} Relationships, belonging, trust
\item \textbf{D (Digital):} Connectivity, access, digital participation
\item \textbf{Eco (Ecological):} Environmental quality, sustainability
\end{itemize}

\textbf{Key properties:}
\begin{itemize}
\item Partially independent (high in one $\neq$ high in others)
\item Systematically interactive (some enable/constrain others)
\item Context-sensitive (each responds to $\Psi$ via $\Gamma$-matrix)
\item Universal (apply across cultures, contexts, time)
\item Measurable (established indicators exist)
\end{itemize}

\textbf{Structural trade-offs revealed by $\Gamma$-matrix:}
\begin{itemize}
\item $\gamma_{S,M} = -0.18$: Market integration can harm social welfare
\item $\gamma_{Eco,E} = -0.19$: Economic development can harm ecological welfare
\item $\gamma_{Eco,M} = -0.21$: Globalization can harm ecological welfare
\end{itemize}

\textbf{Application:}
Each dimension applies to each utility category (INU, KNU, IDN),
creating $3 \times 6 = 18$ base cells of the welfare matrix.
\end{tcolorbox}

\vspace{1em}
\hrule
\vspace{0.5em}
\noindent\textit{Next section: 10.3 The 144-Component Structure}

%%%%%%%%%%%%%%%%%%%%%%%%%%%%%%%%%%%%%%%%%%%%%%%%%%%%%%%%%%%%%%%%%%%%%%%%%%%%%%%
